\lysh{22190_SOLUTION}{1}
ΛΥΣΗ

α) Αρχικά υπολογίζουμε την παράμετρο $p$ της παραβολής. Είναι:
\[ 2p = 1 \Leftrightarrow p = \frac{1}{2} \]
Η εστία $E$ της παραβολής έχει συντεταγμένες
\[ E\left(\frac{p}{2}, 0\right) = \left(\frac{1}{4}, 0\right) \]
Η διευθετούσα $(\delta)$ της παραβολής έχει εξίσωση
\[ x = -\frac{p}{2} = -\frac{1}{4} \]

β) Εξετάζουμε αν οι συντεταγμένες του σημείου $A$ επαληθεύουν την εξίσωση της παραβολής. Για $x=1$ και $y=-1$ είναι:
\[ (-1)^2 = 1 \]
Η τελευταία ισότητα είναι αληθής, οπότε το σημείο $A$ είναι σημείο της παραβολής.

γ) Η εξίσωση της εφαπτομένης της παραβολής $y^2 = 2px$ στο σημείο της $A(x_1,y_1)$ είναι:
\[ yy_1 = p(x + x_1) \]
Αφού δίνεται $A(1,-1)$, η ζητούμενη εξίσωση θα είναι:
\[ -1 \cdot y = \frac{1}{2}(x + 1) \text{ ή } x + 2y + 1 = 0 \]

% TODO: TikZ σχήμα
\endlysh{}